% Options for packages loaded elsewhere
\PassOptionsToPackage{unicode}{hyperref}
\PassOptionsToPackage{hyphens}{url}
\documentclass[
  american,
  11pt,
]{article}
\usepackage{xcolor}
\usepackage[margin=25mm]{geometry}
\usepackage{amsmath,amssymb}
\setcounter{secnumdepth}{-\maxdimen} % remove section numbering
\usepackage{iftex}
\ifPDFTeX
  \usepackage[T1]{fontenc}
  \usepackage[utf8]{inputenc}
  \usepackage{textcomp} % provide euro and other symbols
\else % if luatex or xetex
  \usepackage{unicode-math} % this also loads fontspec
  \defaultfontfeatures{Scale=MatchLowercase}
  \defaultfontfeatures[\rmfamily]{Ligatures=TeX,Scale=1}
\fi
\usepackage{lmodern}
\ifPDFTeX\else
  % xetex/luatex font selection
\fi
% Use upquote if available, for straight quotes in verbatim environments
\IfFileExists{upquote.sty}{\usepackage{upquote}}{}
\IfFileExists{microtype.sty}{% use microtype if available
  \usepackage[]{microtype}
  \UseMicrotypeSet[protrusion]{basicmath} % disable protrusion for tt fonts
}{}
\makeatletter
\@ifundefined{KOMAClassName}{% if non-KOMA class
  \IfFileExists{parskip.sty}{%
    \usepackage{parskip}
  }{% else
    \setlength{\parindent}{0pt}
    \setlength{\parskip}{6pt plus 2pt minus 1pt}}
}{% if KOMA class
  \KOMAoptions{parskip=half}}
\makeatother
% definitions for citeproc citations
\NewDocumentCommand\citeproctext{}{}
\NewDocumentCommand\citeproc{mm}{%
  \begingroup\def\citeproctext{#2}\cite{#1}\endgroup}
\makeatletter
 % allow citations to break across lines
 \let\@cite@ofmt\@firstofone
 % avoid brackets around text for \cite:
 \def\@biblabel#1{}
 \def\@cite#1#2{{#1\if@tempswa , #2\fi}}
\makeatother
\newlength{\cslhangindent}
\setlength{\cslhangindent}{1.5em}
\newlength{\csllabelwidth}
\setlength{\csllabelwidth}{3em}
\newenvironment{CSLReferences}[2] % #1 hanging-indent, #2 entry-spacing
 {\begin{list}{}{%
  \setlength{\itemindent}{0pt}
  \setlength{\leftmargin}{0pt}
  \setlength{\parsep}{0pt}
  % turn on hanging indent if param 1 is 1
  \ifodd #1
   \setlength{\leftmargin}{\cslhangindent}
   \setlength{\itemindent}{-1\cslhangindent}
  \fi
  % set entry spacing
  \setlength{\itemsep}{#2\baselineskip}}}
 {\end{list}}
\usepackage{calc}
\newcommand{\CSLBlock}[1]{\hfill\break\parbox[t]{\linewidth}{\strut\ignorespaces#1\strut}}
\newcommand{\CSLLeftMargin}[1]{\parbox[t]{\csllabelwidth}{\strut#1\strut}}
\newcommand{\CSLRightInline}[1]{\parbox[t]{\linewidth - \csllabelwidth}{\strut#1\strut}}
\newcommand{\CSLIndent}[1]{\hspace{\cslhangindent}#1}
\ifLuaTeX
\usepackage[bidi=basic,shorthands=off]{babel}
\else
\usepackage[bidi=default,shorthands=off]{babel}
\fi
\ifLuaTeX
  \usepackage{selnolig} % disable illegal ligatures
\fi
\setlength{\emergencystretch}{3em} % prevent overfull lines
\providecommand{\tightlist}{%
  \setlength{\itemsep}{0pt}\setlength{\parskip}{0pt}}
\usepackage{bookmark}
\IfFileExists{xurl.sty}{\usepackage{xurl}}{} % add URL line breaks if available
\urlstyle{same}
\hypersetup{
  pdftitle={Heterogeneous initial states: weak instruments and independent-record inference},
  pdflang={en-US},
  hidelinks,
  pdfcreator={LaTeX via pandoc}}

\title{Heterogeneous initial states: weak instruments and
independent-record inference}
\author{}
\date{}

\begin{document}
\maketitle

\subsection{S16. Heterogeneous initial states and independent weak
instruments}\label{s16.-heterogeneous-initial-states-and-independent-weak-instruments}

Consider independent records \(Y_r=x_{\theta_0,u_r}+\eta_r\) of a
deterministic polynomial ODE. The initial states may form a
deterministic triangular array; alternatively they may be iid draws from
a compact set, with the following conditions imposed conditionally on
them. The sensor covariance is \(C_\nu\). Each Hermite-corrected,
exactly integrated weak moment obeys

\[
E\{m_r(\theta_0,\nu_0)\mid u_r\}=0,
\qquad m_r(\theta,\nu)=b_r-G_r(\nu)\theta.
\]

This identity follows by taking expectations of the corrected polynomial
and applying the ODE weak identity separately to each trajectory. Its
validity is uniform over allowed initial states; the distribution of
\(u_r\) need not be estimated. The numerical version needs the
quadrature condition below.

\subsubsection{S16.1. Numbered assumptions and fixed-noise estimating
equations}\label{s16.1.-numbered-assumptions-and-fixed-noise-estimating-equations}

\textbf{(A1)} The coefficient vector has fixed dimension \(p\) and
belongs to an open identifiable neighborhood. Trajectories and their
first coefficient/initial-state sensitivities exist on \([0,T]\).
Initial states lie in a compact set on which the selected moments have
uniformly bounded moments of order \(2+\delta\) for some \(\delta>0\).

\textbf{(A2)} Records and any auxiliary randomizations are independent
across \(r\). Noise is Gaussian with known positive covariance
\(C_{\nu_0}\). Weak tests and quadrature rules are fixed independently
of each record's noise. The coefficient structure is supplied.

\textbf{(A3)} For every record there are two conditionally independent
views \(Y_r^+\) and \(Y_r^-\) with mean \(x_{\theta_0,u_r}\) and
covariance \(2C_{\nu_0}\). An instrument \(L_r^\pm\) is measurable with
respect to its own view and independent auxiliary training data.
Corrected moments \(m_r^\pm\) use the variance of their own view. The
instrument map is fixed conditional on that training data and has the
moment bounds required in (A1).

\textbf{(A4)} Put \(B_r=(L_r^+G_r^-+L_r^-G_r^+)/2\),
\(a_r=(L_r^+b_r^-+L_r^-b_r^+)/2\) and \(\psi_r=a_r-B_r\theta_0\).
Conditional on the training data, \(n^{-1}\sum E B_r\to B\) with \(B\)
nonsingular, and \(n^{-1}\sum\operatorname{Var}(\psi_r)\to V_\psi\)
positive definite. A law of large numbers holds for \(B_r\) and the
Lindeberg condition holds for \(\psi_r\).

\textbf{(A5)} Exact integration is used, or the average instrumented
quadrature defect is \(o(n^{-1/2})\). This condition concerns centered
continuous-parameter inference. A vanishing average defect alone
suffices for consistency.

\textbf{Theorem 3 (heterogeneous-record weak estimating equations).}
Under (A1)--(A5), conditional on the independent training data,

\[
\widehat\theta_n=\left(\sum_{r=1}^n B_r\right)^{-1}\sum_{r=1}^n a_r
\ \xrightarrow{P}\ \theta_0,
\qquad
\sqrt n(\widehat\theta_n-\theta_0)
\ \Rightarrow\ N(0,B^{-1}V_\psi B^{-\top}).
\]

The empirical covariance of the \(p\)-dimensional record scores,
combined with the mean derivative \(\bar B_n\), consistently estimates
this sandwich. Thus a rich weak family requires a \(p\times p\)
empirical score covariance, while its within-record weights can come
from the full model covariance in Proposition 3. Fixed-noise consistency
applies to the stated independent-instrument construction. It does not
follow merely by inserting a same-record pilot into
\(D_r^\top\Omega_r^{-1}\).

\textbf{Proof.} Conditional on \(u_r\) and training data, \(L_r^+\) is
independent of \(m_r^-(\theta_0,\nu_0)\), and conversely. Hence

\[
E(\psi_r\mid u_r,\text{training})
=\tfrac12 E L_r^+\,E m_r^-+
\tfrac12 E L_r^-\,E m_r^+=0.
\]

The exact identity \(\widehat\theta_n-\theta_0=\bar B_n^{-1}\bar\psi_n\)
gives consistency by the law of large numbers, followed by the displayed
CLT by Lindeberg and Slutsky. The empirical score covariance converges
to \(V_\psi\): each conditional mean is zero, so variation of the
deterministic trajectories creates no between-record mean term.
Replacing \(\theta_0\) by the consistent estimate changes this
covariance by \(o_P(1)\) under the moment conditions. The quadrature
assumption makes its centered contribution \(o_P(1)\) after
multiplication by \(\sqrt n\). This is the standard
estimating-equation/GMM argument {[}1, 2{]}, with the
independent-instrument identity addressing the incidental-initial-state
setting. \(\square\)

\subsubsection{S16.2. Gaussian construction and local information
reference}\label{s16.2.-gaussian-construction-and-local-information-reference}

For known \(C_{\nu_0}\), draw \(\zeta_r\sim N(0,C_{\nu_0})\)
independently and set \(Y_r^\pm=Y_r\pm\zeta_r\). Their cross-covariance
is \(C_{\nu_0}-C_{\nu_0}=0\); jointly Gaussian views are independent,
have identical mean and covariance \(2C_{\nu_0}\). Together they recover
the original record. This Gaussian randomization is an established
data-fission construction {[}3{]}. It is used here to check finite-noise
instruments, not proposed as a new general splitting method.

For a supplied clean trajectory, write \(D_r=E G_r\) and
\(\Omega_r=\operatorname{Cov}(m_r(\theta_0,\nu_0))\). The deterministic
optimal instrument within the prescribed linear moments is
\(L_r=D_r^\top\Omega_r^{-1}\). Direct calculation gives the oracle
covariance

\[
\left(\sum_r D_r^\top\Omega_r^{-1}D_r\right)^{-1}.
\]

This is a fixed-family moment reference at finite noise. Proposition 3
evaluates \(\Omega_r\) including all Hermite orders. A noisy
\(G_r^\top\) in place of \(D_r^\top\) introduces \(E(G_r^\top m_r)\) and
can invalidate the zero-mean score. Independent instruments preserve the
mean identity while generally changing the finite-noise covariance from
this oracle expression.

In the local Gaussian experiment, each record has coefficient
sensitivity \(S_r\) and initial sensitivity \(Z_r\). The joint
information for \(\theta,u_1,\ldots,u_n\) has separate nuisance blocks.
Its coefficient Schur complement is exactly

\[
\mathcal I_{e,n}=\sum_{r=1}^n
\left(\mathcal I_{\theta\theta}(u_r)-
\mathcal I_{\theta u}(u_r)\mathcal I_{uu}(u_r)^{-1}
\mathcal I_{u\theta}(u_r)\right)
=\sum_r\mathcal I_e(u_r).
\]

Theorem 2 applies record by record. With known shape and noise
covariance \(\epsilon^2 C\), a fixed test span gives the order-one
moment information \(D_r^\top\Omega_{r,1}^{-1}D_r\). Dense spans
approach \(\mathcal I_e(u_r)\) for any fixed finite collection of
regular trajectories. Thus \((\sum_r\mathcal I_e(u_r))^{-1}\) is the
sharp local reference for boundary-free influences in the product
experiment. This exact block calculation alone supplies neither a
fixed-noise incidental-parameter MLE theorem nor a growing-record,
growing-span uniform efficiency result.

The symmetric independent-view score also preserves the leading
influence under the following additional local conditions: its
instruments converge in \(L^{2+\delta}\) to the deterministic
correct-covariance instruments; their perturbation is \(O_P(\epsilon)\);
weak moments have an \(L^2\) expansion with \(O(\epsilon^2)\) remainder;
and covariance solves remain regular for the fixed span. Then

\[
\tfrac12\{L_r^+m_r^-+L_r^-m_r^+\}
=L_{r,0}\,m_{r,1}(\eta_r)+O_{L^2}(\epsilon^2).
\]

Indeed the first-order moments are linear in \(\eta_r\pm\zeta_r\), their
average cancels \(\zeta_r\), and instrument perturbations multiply
moments of order \(\epsilon\). The factor two in each view's covariance
scales \(L_{r,0}\) uniformly and cancels from the normalized estimating
equation. Combined with Theorem 2, this proves the conditional local
product-experiment efficiency statement for fixed \(n\) and sequential
span enrichment. The fixed-noise CLT above and the local information
reference describe complementary properties of the same weak scores;
joint rates remain open.

\subsubsection{S16.3. Covariance estimation and common-trajectory
sufficiency}\label{s16.3.-covariance-estimation-and-common-trajectory-sufficiency}

For an independently calibrated covariance estimate \(\widehat\nu\),
asymptotic linearity of the averaged score adds \(H\xi\) to its
influence, where \(H=E\partial_\nu\psi\) and \(\xi\) is the calibration
influence with the appropriate relative sample-size scaling. The
Newey--McFadden nuisance argument then gives its sandwich {[}2{]}. For
estimated-covariance Gaussian fission, calibration error also perturbs
cross-view independence; its contribution must be included or be
negligible at the inference scale. The known-covariance construction
proves the stated result without presuming those extra conditions.

Uncorrected residual-profile estimates from separate fixed-horizon
trajectory fits need not be consistent as records accumulate: each fit
estimates a new initial state with nonvanishing error. Even in a linear
Gaussian model, fitting a mean projects part of the error out of the
residual and changes its expected quadratic form. A residual estimator
therefore requires a demonstrated bias correction or a separate
consistent calibration scheme. Pair differences eliminate the mean only
for equal trajectories. These facts prevent transferring the
common-trajectory pair estimator to heterogeneous initial states.

If all records instead share one initial state and known covariance, the
Gaussian path average is sufficient for the mean parameter: the sample
contrasts are parameter-free independent Gaussian observations. Applying
a single-record procedure to the path average reduces covariance to
\(C/n\). For a nonlinear polynomial, a Hermite covariance term of order
\(q\) then scales as \(n^{-q}\), whereas averaging separate corrected
moments retains its fixed-amplitude contribution divided by \(n\). Path
averaging is the natural reference in this common-trajectory experiment.
The old fixed-family GMM calculation is retained as Lemma 2 and a
calibration diagnostic, rather than as an efficiency contribution.

\subsection*{References}\label{references}
\addcontentsline{toc}{subsection}{References}

\protect\phantomsection\label{refs}
\begin{CSLReferences}{0}{0}
\bibitem[\citeproctext]{ref-hansen1982}
\CSLLeftMargin{{[}1{]} }%
\CSLRightInline{Hansen LP. Large sample properties of generalized method
of moments estimators. \emph{Econometrica} 50 (1982) 1029--1054.
\url{https://doi.org/10.2307/1912775}.}

\bibitem[\citeproctext]{ref-newey1994}
\CSLLeftMargin{{[}2{]} }%
\CSLRightInline{Newey WK, McFadden D. Large sample estimation and
hypothesis testing. \emph{Handbook of econometrics} Elsevier 4 (1994)
2111--2245. \url{https://doi.org/10.1016/S1573-4412(05)80005-4}.}

\bibitem[\citeproctext]{ref-leiner2025}
\CSLLeftMargin{{[}3{]} }%
\CSLRightInline{Leiner J, Duan B, Wasserman L, Ramdas A. Data fission:
Splitting a single data point. \emph{Journal of the American Statistical
Association} 120 (2025) 135--146.
\url{https://doi.org/10.1080/01621459.2023.2270748}.}

\end{CSLReferences}

\end{document}
